\chapter{Espacios métricos}
\label{ch:espacios-metricos}
\index{Espacio métrico}\index{Métrica}\index{Distancia}\index{Desigualdad de Cauchy--Schwarz}\index{Métricas equivalentes}\index{Subespacio métrico}

\section{Definición}

Con los números reales ya construidos, el siguiente paso natural en la evolución del rigor es abstraer la noción de distancia. El análisis matemático no estudia únicamente funciones sobre $\mathbb{R}$; sus objetos fundamentales son los espacios métricos, introducidos por Maurice Fréchet en su tesis de 1906 como una generalización que unifica el tratamiento de sucesiones, funciones, conjuntos y transformaciones.

\begin{definition}
\label{def:metrica}
Sea $X$ un conjunto no vacío. Una \emph{métrica} sobre $X$ es una función $d\colon X\times X\to\mathbb{R}$ que satisface, para todos $x,y,z\in X$:
\begin{enumerate}
  \item[(i)] \emph{No negatividad y separación}: $d(x,y)\geq 0$, y $d(x,y)=0$ si y solo si $x=y$;
  \item[(ii)] \emph{Simetría}: $d(x,y)=d(y,x)$;
  \item[(iii)] \emph{Desigualdad triangular}: $d(x,z)\leq d(x,y)+d(y,z)$.
\end{enumerate}
Al par $(X,d)$ se le llama \emph{espacio métrico}.
\end{definition}

La condición (iii) es la que confiere a la métrica su carácter de medida de proximidad compatible con una estructura de límite. Intuitivamente, dice que ir directamente de $x$ a $z$ nunca puede ser más largo que pasar por un intermediario $y$.

\begin{example}
\label{ej:euclidea}
El espacio $\mathbb{R}^n$ dotado de la distancia euclidiana
\begin{equation}\label{eq:euclidea}
  d_2(x,y)=\biggl(\sum_{i=1}^{n}(x_i-y_i)^2\biggr)^{1/2}
\end{equation}
es un espacio métrico. La desigualdad triangular se sigue de la desigualdad de Cauchy--Schwarz.
\end{example}

\begin{example}
\label{ej:discreta}
Sobre cualquier conjunto no vacío $X$, la \emph{métrica discreta}
\[
  d(x,y)=\begin{cases}0&\text{si }x=y,\\1&\text{si }x\neq y\end{cases}
\]
es una métrica. La desigualdad triangular es inmediata: si $x\neq z$, entonces al menos uno de $d(x,y)$ o $d(y,z)$ es $1$.
\end{example}

\begin{example}
\label{ej:supremo}
Sea $C([a,b])$ el conjunto de las funciones continuas $f\colon[a,b]\to\mathbb{R}$. La fórmula
\begin{equation}\label{eq:supremo}
  d(f,g)=\sup_{t\in[a,b]}|f(t)-g(t)|
\end{equation}
define la \emph{métrica del supremo} sobre $C([a,b])$.
\end{example}

\section{La desigualdad de Cauchy--Schwarz}

Antes de continuar con más ejemplos, demostramos la herramienta básica que justifica la desigualdad triangular de la métrica euclidiana.

\begin{theorem}[Cauchy--Schwarz en $\mathbb{R}^n$]
\label{teo:cauchy-schwarz-cap5}
Para cualesquiera $x,y\in\mathbb{R}^n$,
\[
  \biggl|\sum_{i=1}^{n}x_i y_i\biggr|\leq\biggl(\sum_{i=1}^{n}x_i^2\biggr)^{1/2}\biggl(\sum_{i=1}^{n}y_i^2\biggr)^{1/2}.
\]
La igualdad se da si y solo si $x$ e $y$ son proporcionales.
\end{theorem}
\begin{proof}
Si $y=0$ el resultado es trivial. Supongamos $y\neq 0$ y sea $\lambda=\langle x,y\rangle/\langle y,y\rangle$. Entonces
\[
  0\leq\langle x-\lambda y,x-\lambda y\rangle=\|x\|^2-\frac{\langle x,y\rangle^2}{\|y\|^2},
\]
de donde se deduce el teorema.
\end{proof}

\begin{corollary}
\label{cor:euclidiana-metrica}
La función $d_2(x,y)=(\sum|x_i-y_i|^2)^{1/2}$ satisface la desigualdad triangular, luego es una métrica sobre $\mathbb{R}^n$.
\end{corollary}
\begin{proof}
Escribiendo $z_i=x_i-y_i$ y $w_i=y_i-z_i$ (ajustando índices), se aplica Cauchy--Schwarz para acotar $\|x-z\|^2$.
\end{proof}

\section{Ejemplos fundamentales}

\begin{example}[Métricas $\ell^p$ sobre $\mathbb{R}^n$]
\label{ej:lp}
Para $1\leq p<\infty$, la fórmula
\[
  d_p(x,y)=\biggl(\sum_{i=1}^{n}|x_i-y_i|^p\biggr)^{1/p}
\]
define una métrica. El caso $p=2$ es la euclidiana (véase el \cref{cor:euclidiana-metrica}); $p=1$ se llama \emph{taxicab} o \emph{Manhattan}; y el límite $p\to\infty$ produce la métrica del máximo $d_\infty(x,y)=\max_i|x_i-y_i|$.
\end{example}

\begin{example}[Espacio de Hilbert $\ell^2$]
\label{ej:hilbert}
El espacio de las sucesiones cuadrado-sumables,
\[
  \ell^2=\Bigl\{x=(x_n)\in\mathbb{R}^{\mathbb{N}}:\sum_{n=1}^{\infty}x_n^2<\infty\Bigr\},
\]
con la métrica $d(x,y)=(\sum|x_n-y_n|^2)^{1/2}$, es un espacio métrico. La convergencia de la serie que define la métrica se garantiza por la desigualdad de Cauchy--Schwarz en cada truncamiento finito y un paso al límite.
\end{example}

\begin{example}[Métrica del taxi en funciones]
\label{ej:l1}
Sobre $C([a,b])$, la integral
\[
  d_1(f,g)=\int_a^b|f(t)-g(t)|\,dt
\]
define otra métrica, distinta de la del supremo. La convergencia en $d_1$ se llama \emph{convergencia en media}. (La integral es la de Riemann.)
\end{example}

\section{Métricas equivalentes}

\begin{definition}
\label{def:metricas-equivalentes}
Dos métricas $d$ y $d'$ sobre $X$ son \emph{equivalentes} si existen constantes $c,C>0$ tales que
\[
  c\,d(x,y)\leq d'(x,y)\leq C\,d(x,y)\quad\text{para todo }x,y\in X.
\]
\end{definition}

\begin{proposition}
\label{prop:equiv-misma-topologia}
Métricas equivalentes generan la misma familia de conjuntos abiertos (véase el \cref{ch:topologia-metricos}).
\end{proposition}

\begin{example}
En $\mathbb{R}^n$, todas las métricas $\ell^p$ son equivalentes. Sin embargo, la métrica discreta no es equivalente a ninguna de ellas.
\end{example}

\section{Homeomorfismos}

# Homeomorfismos

```latex
\section{Homeomorfismos}

La definición $\epsilon$-$\delta$ de continuidad que se utiliza
habitualmente en cálculo de una variable supone, aunque no siempre se
haga explícito, que estamos trabajando con la métrica euclidiana. En
efecto, una función $f\colon \mathbb{R}\to\mathbb{R}$ es continua en un
punto $a\in\mathbb{R}$ si, para todo $\epsilon>0$, existe
$\delta>0$ tal que, para todo $x\in\mathbb{R}$,
\[
|x-a|<\delta \implies |f(x)-f(a)|<\epsilon.
\]

Esta definición utiliza directamente la distancia euclidiana en
$\mathbb{R}$. La condición $|x-a|<\delta$ describe el intervalo abierto
\[
(a-\delta,a+\delta),
\]
mientras que la condición $|f(x)-f(a)|<\epsilon$ describe el intervalo
abierto
\[
(f(a)-\epsilon,f(a)+\epsilon).
\]
Desde el punto de vista de los espacios métricos, estos intervalos son
precisamente bolas abiertas. Por lo tanto, la idea fundamental de la
continuidad no depende de la estructura particular de $\mathbb{R}$ ni
de la métrica euclidiana, sino que puede formularse en cualquier espacio
métrico mediante sus respectivas bolas abiertas.

\begin{definition}
\label{def:continuidad-metrica}
Sean $(X,d_X)$ y $(Y,d_Y)$ espacios métricos. Decimos que una función
$f\colon X\to Y$ es \emph{continua} en un punto $x_0\in X$ si, para todo
$\epsilon>0$, existe $\delta>0$ tal que, para todo $x\in X$,
\[
d_X(x,x_0)<\delta
\implies
d_Y\bigl(f(x),f(x_0)\bigr)<\epsilon.
\]
\end{definition}

Si una función $f\colon X\to Y$ es continua en todo punto de $X$,
decimos simplemente que $f$ es \emph{continua en $X$} o que es una
\emph{función continua}.

La definición anterior puede interpretarse mediante bolas abiertas. En
efecto, para todo $\epsilon>0$ debe existir $\delta>0$ tal que
\[
f\bigl(B_{d_X}(x_0,\delta)\bigr)
\subseteq
B_{d_Y}\bigl(f(x_0),\epsilon\bigr).
\]
Es decir, los puntos suficientemente cercanos a $x_0$ son enviados por
$f$ a puntos suficientemente cercanos a $f(x_0)$.

\begin{example}
Consideremos la función
\[
f\colon \mathbb{R}\to\mathbb{R},
\qquad
f(x)=2x+1,
\]
donde tanto el dominio como el codominio están dotados de la métrica
euclidiana
\[
d(x,y)=|x-y|.
\]
Demostraremos que $f$ es continua en cualquier punto $x_0\in\mathbb{R}$.

Sea $\epsilon>0$. Necesitamos encontrar $\delta>0$ tal que
\[
|x-x_0|<\delta
\implies
|f(x)-f(x_0)|<\epsilon.
\]
Observamos que
\[
|f(x)-f(x_0)|
=
|(2x+1)-(2x_0+1)|
=
2|x-x_0|.
\]
Por lo tanto, basta elegir
\[
\delta=\frac{\epsilon}{2}.
\]
En efecto, si $|x-x_0|<\delta$, entonces
\[
|f(x)-f(x_0)|
=
2|x-x_0|
<
2\delta
=
\epsilon.
\]
Así, $f$ es continua en $x_0$. Como $x_0$ fue arbitrario, concluimos
que $f$ es continua en todo $\mathbb{R}$.
\end{example}
```


Cuando la continuidad sobre un punto ocurre sobre todo $X$ decimos que la funcion
es continua sobre $X$.
\begin{example}
  Escribir un ejemeplo. 
\end{example}

\section{Subespacios}

\begin{definition}
\label{def:subespacio}
Sea $(X,d)$ un espacio métrico y $Y\subseteq X$ no vacío. La restricción $d|_{Y\times Y}$ define una métrica sobre $Y$, llamada la \emph{métrica inducida}. Al par $(Y,d|_{Y\times Y})$ se le llama \emph{subespacio métrico}.
\end{definition}

\begin{example}
$\mathbb{Q}$, como subconjunto de $\mathbb{R}$ con la métrica usual, es un subespacio métrico. Como veremos en el \cref{ch:completitud}, su comportamiento respecto a sucesiones de Cauchy difiere radicalmente del de $\mathbb{R}$.
\end{example}

\section{Productos de espacios métricos}

\begin{definition}
\label{def:producto-metrico}
Sean $(X,d_X)$ y $(Y,d_Y)$ espacios métricos. Sobre $X\times Y$ se pueden definir varias métricas; la más usuales son:
\begin{itemize}
  \item $d_1((x_1,y_1),(x_2,y_2))=d_X(x_1,x_2)+d_Y(y_1,y_2)$;
  \item $d_2((x_1,y_1),(x_2,y_2))=(d_X(x_1,x_2)^2+d_Y(y_1,y_2)^2)^{1/2}$;
  \item $d_\infty((x_1,y_1),(x_2,y_2))=\max\{d_X(x_1,x_2),d_Y(y_1,y_2)\}$.
\end{itemize}
Todas son equivalentes.
\end{definition}

\section{Espacios métricos inducidos}

\begin{definition}
\label{def:metrica-inducida-topologia}
Sea $(X,d)$ un espacio métrico. Se define la \emph{bola abierta} de centro $x$ y radio $r>0$ como
\[
  B(x,r)=\{y\in X:\,d(x,y)<r\}.
\]
La familia de bolas abiertas genera una topología sobre $X$ (véase el \cref{ch:topologia-metricos}).
\end{definition}

\begin{notation}
\label{not:bolas}
A lo largo del libro, $B(x,r)$ denota la bola abierta, $\overline{B}(x,r)$ la bola cerrada $\{y:d(x,y)\leq r\}$, y $S(x,r)$ la esfera $\{y:d(x,y)=r\}$.
\end{notation}

\begin{exercise}
Demuestre que en un espacio métrico discreto, toda bola abierta $B(x,r)$ con $r\leq 1$ coincide con $\{x\}$.
\end{exercise}

\begin{exercise}
Sea $d$ una métrica sobre $X$. Pruebe que $d'(x,y)=\min\{1,d(x,y)\}$ es también una métrica y que induce la misma topología que $d$.
\end{exercise}

\begin{exercise}
Verifique que la métrica del supremo sobre $C([a,b])$ satisface la desigualdad triangular.
\end{exercise}

\begin{problem}
Sea $(X,d)$ un espacio métrico. Defina $\rho(x,y)=\frac{d(x,y)}{1+d(x,y)}$. Demuestre que $\rho$ es una métrica acotada que genera la misma topología que $d$. ¿Es $\rho$ equivalente a $d$ en el sentido de la \cref{def:metricas-equivalentes}?
\end{problem}
